\documentclass[11pt]{article}
\usepackage[margin=1in]{geometry}
\usepackage{amsmath,amssymb,amsthm}
\usepackage{enumitem}
\usepackage{graphicx}
\usepackage[hidelinks]{hyperref}

\newtheorem{theorem}{Theorem}
\newtheorem{lemma}[theorem]{Lemma}
\newtheorem{proposition}[theorem]{Proposition}
\newtheorem{definition}[theorem]{Definition}
\newtheorem{remark}[theorem]{Remark}

\newcommand{\N}{\mathbb N}
\newcommand{\R}{\mathbb R}
\newcommand{\one}{\mathbf 1}
\newcommand{\omegaone}{\omega_1}
\newcommand{\omegatwo}{\omega_2}
\newcommand{\clspan}{\overline{\operatorname{span}}}
\newcommand{\supp}{\operatorname{supp}}

\title{A ZFC positive answer for overcomplete sets in ordinal interval spaces below \(\omega_2\)}
\author{Agent lane 16}
\date{2026-07-03}

\begin{document}
\maketitle

\noindent\textbf{Status.}
Candidate full solution, likely valid subject to expert review.  This packet
proves a ZFC positive answer to Question 48 from Koszmider's arXiv:2006.00806:
for every ordinal \(\xi<\omega_2\), the Banach space \(C([0,\xi])\) admits an
overcomplete set.

\section{Source question}

The source paper is:

\begin{quote}
P. Koszmider, \emph{On the existence of overcomplete sets in some classical
nonseparable Banach spaces}, arXiv:2006.00806.
\end{quote}

On page 25, Question 48 asks:

\begin{quote}
Can one prove in ZFC that the Banach spaces \(C([0,\xi])\) for all ordinals
\(\xi<\omega_2\) admit overcomplete sets?
\end{quote}

\begin{center}
\includegraphics[width=0.92\linewidth]{figures/open_problem_crop.png}
\end{center}

\section{Definitions and criterion}

\begin{definition}
Let \(X\) be a Banach space.  A set \(D\subseteq X\) is \emph{overcomplete} if
\(|D|=\operatorname{dens}(X)\) and every subset \(A\subseteq D\) with
\(|A|=|D|\) has dense linear span in \(X\).
\end{definition}

We shall use the following elementary criterion repeatedly.

\begin{lemma}\label{lem:functionalcriterion}
Let \(X\) be a Banach space of density \(\omega_1\).  If there is a family
\((x_\alpha)_{\alpha<\omega_1}\subseteq X\) such that
\[
  \{\alpha<\omega_1:x^*(x_\alpha)=0\}
\]
is countable for every nonzero \(x^*\in X^*\), then
\(\{x_\alpha:\alpha<\omega_1\}\) is overcomplete.
\end{lemma}

\begin{proof}
The family has cardinality \(\omega_1\) after deleting possible repetitions;
in the constructions below the points are explicitly distinct.  If an
uncountable subfamily had non-dense linear span, Hahn--Banach would give a
nonzero functional vanishing on it, contradicting the hypothesis.
\end{proof}

We also use the almost coherent injections from Koszmider's proof.  There are
injections \(e_\alpha:\alpha\to\N\) for \(\alpha<\omega_1\), with
\(\N=\{0,1,2,\ldots\}\), such that
\[
   e_\beta\!\upharpoonright\alpha =^* e_\alpha
   \qquad(\alpha<\beta<\omega_1),
\]
and with \(\N\setminus e_\alpha[\alpha]\) infinite.  Consequently, if
\(A\subseteq\omega_1\) is uncountable and \(\gamma<\omega_1\), then there are
an uncountable \(A'\subseteq A\setminus(\gamma+1)\) and an injection
\(g:\gamma+1\to\N\) such that
\[
   e_\alpha\!\upharpoonright(\gamma+1)=g
   \qquad(\alpha\in A').
\]

\section{Proof Intuition}

Koszmider's original construction handles \(C([0,\omega_1])\) by turning the
values of a fixed functional on an uncountable subfamily into zeros of one
analytic function.  The earlier lane-16 packets pushed this through finite
\(\omega_1\)-block sums and then through the first countable block
\(C([0,\omega_1\cdot\omega])\) by assigning disjoint Taylor exponent sets to
the countably many active blocks.

The remaining difficulty is the \(\omega_1^2\)-level.  The interval
\([0,\omega_1^2]\) has a closed skeleton
\[
   B=\{\omega_1\eta:\eta\leq\omega_1\}
\]
homeomorphic to \([0,\omega_1]\), and gaps between consecutive skeleton points
which are copies of the open part of \([0,\omega_1]\).  We split a function
into a skeleton part plus a gap part vanishing on the skeleton.  The skeleton
part uses the strengthened one-block analytic construction.  The gap part uses
endpoint-vanishing block functions, so no endpoint mass has to be detected
inside a gap; endpoint masses belong to the skeleton functional.  For each
selected vector, only countably many gaps are perturbed, and the coherent
injection makes the perturbation norms tend to zero near every limit skeleton
point.  Thus the vector is continuous on \([0,\omega_1^2]\).

For any fixed functional, only countably many gaps and countably many points in
those gaps matter.  One stabilization of the coherent injections freezes all
of them.  The resulting evaluation is an analytic function whose Taylor
coefficients live in pairwise disjoint exponent sets: one set for the skeleton
and one set for each active gap.  Coefficient comparison forces the skeleton
functional and every active gap functional to vanish.  Hence every nonzero
functional has only countably many zeros on the constructed family.

\section{The \texorpdfstring{\(\omega_1^2\)}{omega1 squared} representative}

Let \(L=[0,\omega_1^2]\).  For \(\eta\leq\omega_1\), write
\[
   s_\eta=\omega_1\eta.
\]
The skeleton \(B=\{s_\eta:\eta\leq\omega_1\}\) is homeomorphic to
\([0,\omega_1]\).

There is a norm-one linear extension \(E:C(B)\to C(L)\) defined as follows:
\[
   (Ey)(s_\eta)=y(s_\eta)\quad(\eta\leq\omega_1),
\]
and, for \(s_\eta<t\leq s_{\eta+1}\),
\[
   (Ey)(t)=y(s_{\eta+1}).
\]
This is continuous because if \(\eta_i+1\to\lambda\) in the skeleton, then
\(y(s_{\eta_i+1})\to y(s_\lambda)\); at the right endpoint of each gap the
extension is eventually constant.  At a skeleton point \(s_\lambda\), the
immediate successor \(s_\lambda+1\) gives right-isolating neighbourhoods, so
only the left-hand approach through earlier gaps has to be checked.  Let
\[
   W=\{f\in C(L): f|_B=0\}.
\]
Then \(C(L)=E[C(B)]\oplus W\).

Identify \(C(B)\) with \(C([0,\omega_1])\).  Put
\[
   u_\alpha=\one_{[0,\alpha]}\in C([0,\omega_1])\qquad(\alpha<\omega_1),
\]
and
\[
   Z=\{z\in C([0,\omega_1]):z(0)=z(\omega_1)=0\}.
\]
For \(\eta<\omega_1\), let \(\iota_\eta:Z\to W\) be the insertion into the
closed gap \([s_\eta,s_{\eta+1}]\), using the order homeomorphism with
\([0,\omega_1]\), and zero outside that gap.

\begin{theorem}\label{thm:omega1square}
In ZFC, \(C([0,\omega_1^2])\) admits an overcomplete set.
\end{theorem}

\begin{proof}
Choose pairwise disjoint infinite sets
\[
   P_\infty,P_0,P_1,\ldots\subseteq\N
\]
and enumerations \(p_j:\N\to P_j\), \(j\in\{\infty,0,1,\ldots\}\), such that
\[
   p_j(k)\geq 2k+1\qquad(k<\omega).
\]
Set
\[
   q_\infty(t)=\sum_{k=0}^\infty 2^k t^{p_\infty(k)}
   \qquad(|t|<1/\sqrt2).
\]
Fix \(0<b<1/2\), choose distinct \(r_\alpha\in(b/2,b)\) for
\(\alpha<\omega_1\), and choose positive scalars \(a_m\downarrow0\).

For \(\alpha<\omega_1\), define the skeleton vector \(y_\alpha\in C(B)\cong
C([0,\omega_1])\) by
\[
 y_\alpha
   = r_\alpha^{p_\infty(0)}u_\alpha
     +\sum_{\beta<\alpha}
        r_\alpha^{p_\infty(e_\alpha(\beta)+1)}(u_\beta-u_\alpha)
     +q_\infty(r_\alpha)\one .
\]
For \(m<\omega\), define the endpoint-vanishing gap vector
\[
 z_\alpha^{(m)}
   =\sum_{\beta<\alpha}
        r_\alpha^{p_m(e_\alpha(\beta)+1)}(u_\beta-u_\alpha)
     \in Z.
\]
The series converge absolutely because \(e_\alpha\) is injective and
\(\sum_k b^{p_j(k)}<\infty\).

Now put
\[
   x_\alpha
   =E y_\alpha+
     \sum_{\eta<\alpha}
       a_{e_\alpha(\eta)}\,\iota_\eta\bigl(z_\alpha^{(e_\alpha(\eta))}\bigr)
       \qquad(\alpha<\omega_1).
\]
The displayed sum is well-defined pointwise because each point of
\(L\setminus B\) lies in one gap.  It is continuous: at successor skeleton
points the inserted gap functions vanish at both endpoints; at a limit
skeleton point \(s_\lambda\), the finite sets
\(\{\eta<\alpha:e_\alpha(\eta)\leq M\}\) are bounded below \(\lambda\), so
\(a_{e_\alpha(\eta)}\to0\) as \(\eta\to\lambda\).  The same argument handles
the top point \(s_{\omega_1}\), while the right side of each skeleton point is
isolated by its immediate successor.  Thus \(x_\alpha\in C(L)\).  Also the
\(x_\alpha\)'s are distinct, since
\[
   x_\alpha(s_{\omega_1})=q_\infty(r_\alpha),
\]
and \(q_\infty\) is strictly increasing on \((0,b)\).

We prove the functional separation property.  Let \(\mu\in C(L)^*\) and
suppose that
\[
   \mu(x_\alpha)=0\qquad(\alpha\in A)
\]
for some uncountable \(A\subseteq\omega_1\).  Since \(L\) is scattered, \(\mu\)
is an atomic measure with countable support.

Let \(\nu=E^*\mu\in C(B)^*\).  Let \(S\subseteq\omega_1\) be the countable set
of gaps on which the restriction of \(\mu\) to the open gap
\((s_\eta,s_{\eta+1})\) is nonzero.  For \(\eta\in S\), view this restriction
as a measure \(\mu_\eta\) on \([0,\omega_1]\) with no mass at \(0\) or
\(\omega_1\); it represents the induced functional on \(Z\).

Choose \(\gamma<\omega_1\) large enough so that:
\begin{itemize}[leftmargin=2em]
\item \(S\subseteq[0,\gamma]\);
\item the support of \(\nu\) below the top point of \(B\) is contained in
      \([0,\gamma]\);
\item for every \(\eta\in S\), the support of \(\mu_\eta\) in
      \([0,\omega_1)\) is contained in \([0,\gamma]\).
\end{itemize}
Stabilize the coherent injections on \(\gamma+1\): pass to an uncountable
\(A'\subseteq A\setminus(\gamma+1)\) and an injection \(g:\gamma+1\to\N\) such
that \(e_\alpha\!\upharpoonright(\gamma+1)=g\) for \(\alpha\in A'\).

For the skeleton functional \(\nu\), write
\[
   b_\beta=\nu(u_\beta)\quad(\beta\leq\gamma),\qquad
   d=\nu(\one),
\]
and let \(c\) be the eventual value of \(\nu(u_\xi)\) for \(\xi>\gamma\).  For
each \(\eta\in S\), write
\[
   b_{\eta,\beta}=\mu_\eta(u_\beta)\quad(\beta\leq\gamma),
\]
and let \(c_\eta\) be the eventual value of \(\mu_\eta(u_\xi)\) for
\(\xi>\gamma\).

For \(\alpha\in A'\), with \(t=r_\alpha\), the equality \(\mu(x_\alpha)=0\)
becomes \(F(t)=0\), where
\[
\begin{aligned}
F(t)=&
 c\,t^{p_\infty(0)}
 +\sum_{\beta\leq\gamma}(b_\beta-c)t^{p_\infty(g(\beta)+1)}
 +d\,q_\infty(t)\\
&+\sum_{\eta\in S} a_{g(\eta)}
   \sum_{\beta\leq\gamma}
      (b_{\eta,\beta}-c_\eta)t^{p_{g(\eta)}(g(\beta)+1)} .
\end{aligned}
\]
This is analytic on \(|t|<1/\sqrt2\).  The last sum converges normally on
compact subdiscs because \(\sum_{\eta\in S}\|\mu_\eta\|\leq\|\mu\|\) and
\(\sum_k |t|^{p_m(k)}\) is bounded uniformly in \(m\) on such subdiscs.
The set \(\{r_\alpha:\alpha\in A'\}\) has an accumulation point in
\((b/2,b)\), hence \(F\equiv0\).

Compare coefficients.  The sets \(P_\infty,P_0,P_1,\ldots\) are disjoint.
On \(P_\infty\), the same argument as in the one-block construction applies.
Since the range of \(g\) has infinite complement, coefficients outside
\(\{0\}\cup\{g(\beta)+1:\beta\leq\gamma\}\) give \(d=0\).  The coefficient at
\(p_\infty(0)\) gives \(c=0\), and the coefficients at
\(p_\infty(g(\beta)+1)\) give \(b_\beta=0\) for all \(\beta\leq\gamma\).
Thus \(\nu\) annihilates \(\one\) and all \(u_\xi\)'s; these span a dense
subalgebra of \(C(B)\), so \(\nu=0\).

Fix \(\eta\in S\).  On the exponent set \(P_{g(\eta)}\), only the
\(\eta\)-gap contributes.  Since \(a_{g(\eta)}>0\), the coefficients at
\(p_{g(\eta)}(g(\beta)+1)\) give
\[
   b_{\eta,\beta}=c_\eta\qquad(\beta\leq\gamma).
\]
For \(\beta=0\), \(b_{\eta,0}=0\), because \(\mu_\eta\) has no mass at the
left endpoint of the model interval.  Hence \(c_\eta=0\), and then all
\(b_{\eta,\beta}=0\).  Therefore \(\mu_\eta(u_\xi)=0\) for every
\(\xi<\omega_1\).  Since \(\mu_\eta\) has no endpoint masses, it annihilates
the dense span of the endpoint-vanishing differences \(u_\beta-u_\alpha\) in
the quotient \(Z\); equivalently, the induced functional on \(Z\) is zero.

We have shown that \(\mu\) annihilates \(E[C(B)]\) and every gap summand in
\(W\).  Since \(C(L)=E[C(B)]\oplus W\), it follows that \(\mu=0\).  Thus no
nonzero functional vanishes on uncountably many \(x_\alpha\)'s.  By
Lemma~\ref{lem:functionalcriterion}, \(C([0,\omega_1^2])\) admits an
overcomplete set.
\end{proof}

\section{Classification reduction below \texorpdfstring{\(\omega_2\)}{omega2}}

We now combine the representative cases.  We use the classical isomorphic
classification of ordinal interval \(C(K)\)-spaces due independently to
Gul'ko--Os'kin and Kislyakov.  In the form needed here, for ordinals of
cardinality \(\omega_1\):
\begin{itemize}[leftmargin=2em]
\item finite multiplicities \([0,\omega_1 n+\delta]\), \(n<\omega\) and
      \(\delta<\omega_1\), reduce to finite \(\ell_\infty\)-sums of blocks of
      type at most \([0,\omega_1]\);
\item all countably infinite multiplicities
      \(\omega_1\cdot\rho\leq \xi<\omega_1\cdot(\rho+1)\) with
      \(\omega\leq\rho<\omega_1\) are isomorphic to the representative
      \(C([0,\omega_1\cdot\omega])\);
\item all \(\xi\) with \(\omega_1^2\leq\xi<\omega_2\) are isomorphic to the
      representative \(C([0,\omega_1^2])\).
\end{itemize}
This is the \(\lambda=\omega_1\) instance of Kislyakov's classification; the
same trichotomy is summarized in later ordinal-classification literature.

\begin{theorem}\label{thm:full}
In ZFC, for every ordinal \(\xi<\omega_2\), the Banach space
\(C([0,\xi])\) admits an overcomplete set.
\end{theorem}

\begin{proof}
If \(\xi<\omega_1\), then \(C([0,\xi])\) is separable, so Klee's theorem gives
an overcomplete sequence.

Let \(\omega_1\leq\xi<\omega_2\).  Write
\[
   \xi=\omega_1\cdot\rho+\delta,\qquad \delta<\omega_1.
\]
If \(\rho<\omega\), the earlier finite-block argument applies: \([0,\xi]\)
splits into finitely many clopen ordinal blocks of type at most
\([0,\omega_1]\), and Koszmider's \(C([0,\omega_1])\) criterion plus finite
\(\ell_\infty\)-sum stability give an overcomplete set.

If \(\omega\leq\rho<\omega_1\), the classification quoted above gives
\[
   C([0,\xi])\sim C([0,\omega_1\cdot\omega]).
\]
The countable-block packet proves that \(C([0,\omega_1\cdot\omega])\) admits
an overcomplete set; equivalently, one may use the same sparse-exponent proof
for the scalar plus countable \(c_0\)-sum model.  Since overcompleteness is
preserved by Banach-space isomorphism, \(C([0,\xi])\) admits an overcomplete
set.

Finally, if \(\omega_1\leq\rho\), then \(\xi\geq\omega_1^2\), and the same
classification gives
\[
   C([0,\xi])\sim C([0,\omega_1^2]).
\]
By Theorem~\ref{thm:omega1square}, \(C([0,\omega_1^2])\) admits an
overcomplete set.  Therefore \(C([0,\xi])\) admits an overcomplete set.
\end{proof}

\section{Verification and novelty check}

Local indexes were searched for arXiv:2006.00806, Question 48, ordinal
intervals, \(C([0,\xi])\), \(C([0,\omega_1\cdot\omega])\), and
\(C([0,\omega_1^2])\).  Existing lane-16 packets covered
\(\xi<\omega_1\cdot\omega\) and then \(C([0,\omega_1\cdot\omega])\), but not
the \(\omega_1^2\)-representative.  Fresh literature searches for the exact
question and close phrases did not locate a later paper explicitly resolving
Koszmider's Question 48.

The proof relies on classical ordinal \(C(K)\)-space classification and on the
almost coherent injections already used in Koszmider's source paper.  The new
mathematical step is Theorem~\ref{thm:omega1square}, especially the skeleton
extension plus endpoint-vanishing gap perturbations.

\section{Human review recommendation}

Review as a candidate full solution.  The main checks are:
\begin{enumerate}[leftmargin=2em]
\item continuity of the constructed \(x_\alpha\)'s at limit skeleton points in
      \([0,\omega_1^2]\);
\item the decomposition \(C([0,\omega_1^2])=E[C(B)]\oplus W\);
\item the coefficient comparison on \(P_\infty\) and on the active gap sets
      \(P_{g(\eta)}\);
\item the use of Kislyakov's classification to reduce all
      \(\xi<\omega_2\) to the finite, countable, and \(\omega_1^2\)
      representatives.
\end{enumerate}

\begin{thebibliography}{9}

\bibitem{Koszmider}
P. Koszmider,
\emph{On the existence of overcomplete sets in some classical nonseparable
Banach spaces},
J. Funct. Anal. \textbf{281} (2021), no. 9, 109172;
\href{https://arxiv.org/abs/2006.00806}{arXiv:2006.00806}.

\bibitem{RussoSomaglia}
T. Russo and J. Somaglia,
\emph{Overcomplete sets in non-separable Banach spaces},
Proc. Amer. Math. Soc. \textbf{149} (2021), 701--714;
\href{https://arxiv.org/abs/1912.08690}{arXiv:1912.08690}.

\bibitem{GulkoOskin}
S. P. Gul'ko and A. V. Os'kin,
\emph{Isomorphic classification of spaces of continuous functions on totally
ordered compact sets},
Funct. Anal. Appl. \textbf{9} (1975), 56--57.

\bibitem{Kislyakov}
S. V. Kislyakov,
\emph{Classification of spaces of continuous functions of ordinals},
Siberian Math. J. \textbf{16} (1975), 226--231.

\bibitem{GenzeGulkoKhmyleva}
L. V. Genze, S. P. Gul'ko, and T. E. Khmyleva,
\emph{Classification of spaces of continuous functions on ordinals},
\href{https://arxiv.org/abs/1806.09015}{arXiv:1806.09015}.

\end{thebibliography}

\end{document}
